\subsection{\texorpdfstring{\(\mathcal{L}(\mathbf{E};\mathbf{F})\) 上的 S-有界结构}{\textbackslash mathcal\{L\}(\textbackslash mathbf\{E\};\textbackslash mathbf\{F\}) 上的 S-有界结构}}

设 E 和 F 是两个局部凸空间，S 是 E 的有界子集的一个族。说 \(\mathcal{L}(E; F)\) 的子集 H 对 S-拓扑有界，意思是对每个 \(\mathbf{M} \in \mathfrak{S}\)，0 的每个邻域 \(V\)（在 \(F\) 中）吸收集 \(\mathrm{H(M)} = \bigcup_{u \in H} u(\mathrm{M})\)；这等于说，对每个 \(\mathbf{M} \in \mathfrak{S}\)，集 \(\mathrm{H(M)}\) 在 \(F\) 中有界。等价地，这意味着对 0 的每个邻域 \(V\)（在 \(F\) 中），集 \(\bigcap_{u \in H} u^{-1}(V)\)

吸收 S 的每个子集 M。

命题 9. --- 设 E 和 F 是两个局部凸空间，S 是 E 的有界子集的一个族。则 \(\mathcal{L}(\mathrm{E};\mathrm{F})\) 的每个等度连续子集对 S-拓扑有界。

因为若 \(\mathbf{H}\) 是 \(\mathcal{L}(\mathrm{E};\mathrm{F})\) 的一个等度连续子集，且 \(\mathbf{V}\) 是 \(\mathbf{F}\) 中 0 的一个邻域，则集 \(\bigcap_{u\in \mathbf{H}}u^{-1}(\mathbf{V})\) 是 \(\mathbf{E}\) 中 0 的一个邻域，故吸收 \(\mathbf{E}\) 的每个有界子集。

\(\mathcal{L}(\mathrm{E};\mathrm{F})\) 中对某个 S-拓扑有界的子集未必等度连续，即使 S 是覆盖的且 S 是 E 上的典范有界结构（IV, 第 50 页, 习题 17）。在下一段中我们将以桶型空间之名研究这样的空间 E：\(\mathcal{L}(\mathrm{E};\mathrm{F})\) 的每个简单有界子集都是等度连续的。目前请注意下面的结果：

命题 10. --- 设 E 是一个有界型空间（特别地，是可度量化的局部凸空间），F 是一个局部凸空间。\(\mathcal{L}(\mathrm{E};\mathrm{F})\) 中对有界收敛拓扑有界的每个子集 H 都是等度连续的。

对 F 中 0 的每个凸均衡邻域 V，集 \(\bigcap_{u\in\mathbf{H}}u^{-1}(\mathbf{V})\) 吸收 E 的每个有界子集，故是 E 中 0 的一个邻域；这就证明了 H 等度连续。

